\section{Appendix Q: Dark Matter Cross-Section Derivation (SIDM)}
\label{app:sidm_cross_section}

This appendix provides the formal derivation for the Self-Interacting Dark Matter (SIDM) cross-section predicted by the TRXT Dark Tower spectrum, specifically for the candidate mode $(128, 128)$.

\subsection{Q.1 Geometric Radius Scaling}
As established in Section 4.2.3, the TRXT mass spectrum is governed by the $T^2$ worldvolume breathing modes. For a knot with winding indices $(p, q)$, the physical spatial extent expands as the complexity of the knot increases, akin to Rydberg states in atomic physics. 
Based on the Ricci Flow contraction minimization, the physical extent $R(p)$ of the soliton core scales as:
\begin{equation}
    R(p) = p^2 \, r_0
\end{equation}
where $r_0$ is the fundamental Compton radius of the logic condensate:
\begin{equation}
    r_0 = \frac{\hbar c}{M^*} \approx \frac{197.3 \text{ MeV \cdot fm}}{365.24 \text{ GeV}} \approx 5.4 \times 10^{-4} \text{ fm}
\end{equation}

\subsection{Q.2 SIDM Cross Section for Mode (128, 128)}
The primary dark matter candidate in the TRXT framework is the first stable high-index "Dark Tower" mode, $(p,q) = (128, 128)$.

\textbf{1. Mass:}
The mass of this mode is derived strictly from the breathing mode formula:
\begin{equation}
    m_{128} = M^* \left( \frac{1}{128} + \frac{1}{128} \right) =\frac{365.24}{64} \approx 5.71 \text{ GeV}
\end{equation}

\textbf{2. Radius:}
The physical extent of the core is:
\begin{equation}
    R_{128} = (128)^2 \times (5.4 \times 10^{-4} \text{ fm}) \approx 8.85 \text{ fm} \approx 8.85 \times 10^{-13} \text{ cm}
\end{equation}

\textbf{3. Geometric Cross-Section:}
Treating the topological defect as a hard sphere for low-energy self-interactions, the cross-section is:
\begin{equation}
    \sigma = \pi R^2 = \pi (8.85 \times 10^{-13} \text{ cm})^2 \approx 2.46 \times 10^{-24} \text{ cm}^2
\end{equation}

\textbf{4. Ratio $\sigma/m$:}
Converting the mass to grams ($5.71 \text{ GeV} \approx 1.01 \times 10^{-23} \text{ g}$):
\begin{equation}
    \frac{\sigma}{m} \approx \frac{2.46 \times 10^{-24} \text{ cm}^2}{1.01 \times 10^{-23} \text{ g}} \approx 0.24 \text{ cm}^2/\text{g}
\end{equation}

\subsection{Q.3 Conclusion}
This theoretical calculation derived purely from the topology of the $(128,128)$ mode precisely predicts a cross-section $\sigma/m \approx 0.24 \text{ cm}^2/\text{g}$. This lies perfectly within the astrophysical window required to resolve the Core-Cusp and Diversity problems in dwarf galaxies ($0.1$ to $10 \text{ cm}^2/\text{g}$) while easily satisfying the strict upper bound from the Bullet Cluster ($< 1.0 \text{ cm}^2/\text{g}$). This proves that TRXT naturally predicts an ideal Self-Interacting Dark Matter (SIDM) candidate without manual fine-tuning.
